\subsection{Why Did Language Evolve?}
\label{sec:drivers}
\slidesrc{Intro: 13--18}

Knowing \emph{when} language appeared does not tell us \emph{why}. This part of the
course follows the linguist Sverker Johansson (\emph{The Dawn of Language: How We Came
to Talk}), for whom one major driver was the \textbf{obstetrical dilemma}.

\subsubsection{The obstetrical dilemma}
\slidesrc{Intro: 13--16}

Two evolutionary trends pulled the human body in opposite directions:
\begin{itemize}[nosep]
  \item \textbf{Upright walking (bipedalism)} restructured the human pelvis into a
        basket-like shape that supports the internal organs and makes walking efficient,
        but \emph{narrows the birth canal}.
  \item \textbf{Large brains}: hominid brain size increased dramatically over millions
        of years, so the baby's head became large compared with the mother's pelvis.
\end{itemize}
\textbf{Consequence}: labour and delivery are difficult and dangerous in the human
species (not only in humans, but especially in humans).

Evolution's ``compromise'' was to have babies born \textbf{immature}: their brain keeps
growing a lot after birth. Human newborns cannot even cling to their mother, as the
babies of all other higher primates do, and they need the protection of their family
or social group for a very long time.

\begin{nbintuition}[how this leads to language]
The course asks: how might these facts relate to the birth of language? Follow the
chain:
\begin{enumerate}[nosep]
  \item A difficult birth means a mother needs help: \textbf{midwives} and social
        support. Giving birth becomes a social, cooperative event.
  \item A baby that cannot cling must be carried or \textbf{put down}. A mother who puts
        the baby down to gather food can still soothe and monitor it from a distance with
        her \textbf{voice} (``motherese'').
  \item A long childhood is a long \textbf{window for learning}: there is time to
        absorb something as complex as a language from the community.
\end{enumerate}
\end{nbintuition}

Table~\ref{tab:pressures} summarises the argument.

\begin{center}
\small
\begin{tabularx}{\linewidth}{@{}>{\raggedright\arraybackslash}p{2.6cm}
                               >{\raggedright\arraybackslash}p{4.1cm} X@{}}
\toprule
\textbf{Evolutionary pressure} & \textbf{Physical consequence} & \textbf{Behavioural and language outcome} \\
\midrule
Bipedalism and big brains & Narrow birth canal; difficult, dangerous labour &
  Need for \textbf{midwives and social support}, forcing cooperative social structures \\
Neurological immaturity & Babies are born completely helpless and cannot cling &
  Mothers must \textbf{``put the baby down''}, using vocalisations to soothe and monitor
  them from afar \\
Extended childhood & Years of dependency requiring group protection &
  Long childhoods provide a \textbf{massive window for learning} complex language from
  the community \\
\bottomrule
\end{tabularx}
\captionof{table}{From the obstetrical dilemma to language.}
\label{tab:pressures}
\end{center}

\subsubsection{Other key drivers}
\slidesrc{Intro: 17}

Besides the obstetrical dilemma, three more drivers are usually cited:
\begin{itemize}[nosep]
  \item \textbf{Survival and cooperation}: early humans needed to coordinate group hunts,
        plan defences against predators and share crucial resources.
  \item \textbf{Social bonding}: communication grew from basic physical grooming into
        storytelling around fires, building community trust and deep empathy (this is
        developed in Section~\ref{sec:gossip}).
  \item \textbf{Abstract thought}: language enabled humans to refer to things not
        immediately present, like the past, the future or imaginary cultural concepts.
        This is displacement again (Section~\ref{sec:human-vs-animal}).
\end{itemize}

\subsubsection{Three theories of how language emerged}
\slidesrc{Intro: 18}

There are three main families of theories. They differ on \emph{speed} (gradual or
sudden) and on the \emph{starting medium} (voice or gestures).

\begin{center}
\small
\begin{tabularx}{\linewidth}{@{}>{\raggedright\arraybackslash}p{1.9cm} *{3}{>{\raggedright\arraybackslash}X}@{}}
\toprule
 & \textbf{Continuity theory} & \textbf{Discontinuity theory} & \textbf{Gesture-first theory} \\
\midrule
Core idea & Language evolved \textbf{gradually} over millions of years out of animal
  communication systems &
  Language appeared \textbf{suddenly} as a unique, exclusively human genetic mutation &
  Language originated from \textbf{physical gestures} and manual signs before shifting
  to speech \\[2pt]
Key mechanism & Natural selection slowly shaped the vocal tract and brain tissue to
  handle complex communication &
  A sudden cognitive ``leap'' or single genetic mutation re-wired the brain for
  universal grammar &
  Bipedalism freed the hands, allowing early humans to use detailed hand and facial
  signs to communicate \\[2pt]
Primary driver & The gradual need for social bonding and survival coordination in
  expanding groups &
  An internal computational upgrade in the brain, prioritising thought organisation over
  communication &
  The mechanical advantage of visual ``miming'' before vocalisation evolved to take
  over \\[2pt]
Famous proponents & Charles Darwin & Noam Chomsky & Michael Corballis \\
\bottomrule
\end{tabularx}
\captionof{table}{The three main theories on the emergence of language.}
\label{tab:theories}
\end{center}

\begin{nbintuition}[telling them apart]
\begin{itemize}[nosep]
  \item \textbf{Darwin}: language is like any other trait, shaped slowly by natural
        selection; ``more of'' what animals already have.
  \item \textbf{Chomsky}: language is primarily a tool for \emph{thinking}, and its core
        (the ability to combine units recursively, ``universal grammar'') appeared in one
        jump; communication came second.
  \item \textbf{Corballis}: language started in the \emph{hands}; the voice took over
        later because it works in the dark, at a distance and with the hands busy.
\end{itemize}
Chomsky's name will come back in this course for a very different reason: his work on
formal grammars underlies the syntactic parsing we will study later.
\end{nbintuition}
